% Generated draft: requires mathematical and source review.
\documentclass[11pt]{article}
\usepackage[margin=25mm]{geometry}
\usepackage{fontspec}
\setmainfont{Latin Modern Roman}
\usepackage{amsmath,amssymb,mathtools,unicode-math}
\setmathfont{Latin Modern Math}
\usepackage{enumitem,needspace,xurl,hyperref}
\setlength{\parindent}{0pt}
\setlength{\parskip}{5pt}
\setlength{\emergencystretch}{3em}
\newcommand{\N}{\mathbb{N}}
\newcommand{\Z}{\mathbb{Z}}
\newcommand{\Q}{\mathbb{Q}}
\newcommand{\R}{\mathbb{R}}
\newcommand{\C}{\mathbb{C}}
\newcommand{\F}{\mathbb{F}}
\newcommand{\A}{\mathbb{A}}
\newcommand{\PP}{\mathbb{P}}
\newcommand{\E}{\mathbb{E}}
\newcommand{\eps}{\varepsilon}
\newcommand{\dd}{\,\mathrm{d}}
\begin{document}
\section*{UnsolvedMath ID: 700495}
\textbf{Source problem code:} catalogue entry
\subsection{Definitions and mathematical statement}
Let \(\mathfrak X\) and \(\mathfrak Y\) be classes of groups, and let \(\mathfrak X\cap\mathfrak Y\) denote the class of groups belonging to both. Using the quotient convention, a group \(G\) is \emph{residually a \(\mathfrak C\)-group} if
\[
\forall g\in G\setminus\{1\}\;\exists N\trianglelefteq G
\quad\bigl(g\notin N\ \text{and}\ G/N\in\mathfrak C\bigr).
\]
Thus each nonidentity element must survive in some quotient belonging to \(\mathfrak C\); the quotient may depend on the element.

The supplied problem asks under what conditions the implication
\[
\left\{
\begin{aligned}
&G\text{ is residually a }\mathfrak X\text{-group},\\
&G\text{ is residually a }\mathfrak Y\text{-group}
\end{aligned}
\right.
\quad\Longrightarrow\quad
G\text{ is residually a }(\mathfrak X\cap\mathfrak Y)\text{-group}
\]
holds. Explicitly, its conclusion requires
\[
\forall g\in G\setminus\{1\}\;\exists N\trianglelefteq G
\quad\bigl(g\notin N\ \text{and}\ G/N\in\mathfrak X\cap\mathfrak Y\bigr),
\]
whereas the hypotheses allow separate quotients in \(\mathfrak X\) and in \(\mathfrak Y\).

A class-level interpretation is to characterize the pairs \((\mathfrak X,\mathfrak Y)\) for which this implication holds for every group \(G\). The source does not specify whether the requested conditions concern only these classes, the group \(G\), or both. It supplies no closure assumptions on the classes and no explicit convention for residuality. Consequently, the displayed formulation uses a stated convention and preserves the request for conditions; it does not assert that the implication holds for arbitrary classes and groups.
\subsection{Short English statement}
Determine when being residually in each of two classes of groups guarantees being residually in their intersection. Under the quotient convention, every nonidentity element must survive in a quotient belonging to both classes. The source leaves the permitted classes and the scope of the requested conditions unspecified.
\subsection{Sources}
\begin{itemize}
\item Local UnsolvedMath registry, supplied complete record: permanent ID \(700495\), source code \texttt{catalogue entry}, title \emph{Kourovka Notebook Problem 8.45}, upstream revision \texttt{d371ee917ac9647fbb06cd6fc529deec6f3dda82}.
\item Attributed in the supplied record to Yu. I. Merzlyakov, proposed in 1982; original problem number 8.45 in \emph{The Kourovka Notebook (Unsolved Problems in Group Theory)}, 8th issue (1982). The supplied background also identifies the 21st issue edition (2026), \texttt{arXiv:1401.0300v48}. These bibliographic details are supplied attribution and have not been independently verified.
\item The supplied collection description credits E. I. Khukhro and V. D. Mazurov as editors of the collection covering Issues 1--21 (1965--2026). The registry labels this problem \emph{open}; this is a supplied status label, not an independently verified literature assessment.
\end{itemize}
\end{document}
